\section{Limitations}
\label{sec:limitations}

\begin{enumerate}
    \item \textbf{Single machine and material.}
    All experiments use a single Bantam Tools desktop CNC mill (air-cut runs and dry UHMW-PE active cuts).
    Industrial CNC machines vary widely in kinematics, controllers, and power levels.
    Metals produce different sensor signatures than UHMW-PE.
    The reconstruction-mismatch principle should generalize, but the specific AUROCs reported here require empirical validation on other platforms and materials.

    \item \textbf{Synthetic, non-adaptive attacks; integrity-check threat model.}
    All attacks are synthetically generated by modifying G-code tokens post-alignment.
    The headline AUROCs assume a \emph{non-adaptive} attacker.
    A first-order adaptive attacker that selects, per window, the (still sabotage-relevant) perturbation minimizing the detector score drops single-score detection to near chance (AUROC = 0.54; Section~\ref{sec:results_revision}).
    Robust detection against such an adversary is unsolved.
    The evaluation further assumes the sensors reflect the \emph{original} program while the claimed G-code is the attacker's.
    This assumption fits post-transmission/at-rest integrity verification against an authenticated original-run trace.
    It does \emph{not} cover attacks that modify G-code \emph{before} execution so that the sensors reflect the tampered process (the canonical dr0wned/Sturm scenario), a regime this method cannot detect.

    \item \textbf{Small dataset.}
    The 120-file dataset (${\sim}545$ sensor windows) limits statistical power.
    Several attack-type/method combinations have confidence intervals encompassing chance (e.g., V7 $S_{\text{mean}}$ on A2 at $\delta = 0.001$~in: AUROC = $0.697 \pm 0.151$).
    The damage-condition classes (4--5 files each) are particularly underrepresented.

    \item \textbf{Known toolpath vocabulary.}
    The system is trained on a fixed set of 9 toolpath classes with a 712-token vocabulary.
    Novel G-code programs using different commands or coordinate ranges may trigger false positives.
    Introducing new programs therefore requires periodic retraining.

    \item \textbf{Sensor trustworthiness assumption.}
    The threat model assumes sensor signals cannot be spoofed (Section~\ref{sec:adversary}).
    An attacker who simultaneously modifies G-code and injects false sensor data (via physical tampering, bus interception, or software compromise) could evade detection entirely.
    The assumption holds for hardwired sensors with shielded cables and dedicated data paths, but breaks for wireless or network-accessible sensor interfaces.
    Mitigation requires hardware-level tamper protection (secure enclaves, hardware attestation) for the sensor subsystem.

    \item \textbf{Offline evaluation only.}
    All experiments run in batch mode.
    Real-time deployment introduces latency constraints, windowing edge effects, online threshold adaptation, and CNC controller integration, none of which are evaluated here.
    The measured 5.8~ms/window (GPU) decoder inference latency suggests computational feasibility.

    \item \textbf{Binary detection scope.}
    The system produces a binary normal/anomalous decision without diagnosing attack type, identifying modified parameters, or recommending corrective action.
    Attack diagnosis would require labeled attack examples or interpretability methods tracing anomaly scores to specific token positions.

    \item \textbf{Supervised anomaly classifier.}
    The XGBoost anomaly classifier (Section~\ref{sec:results_anomaly_head}) uses supervised training with labeled attacks.
    Its applicability is therefore limited to attack types seen during training.
\end{enumerate}
